\section{Introduction}\label{sec:intro}
A trader who has already committed to sell a fixed quantity faces a different problem from an investor deciding whether to hold an asset. The execution trader allocates the required quantity across available times. Predicting the final price can help with that task if it also reveals how prices are expected to evolve before completion. By itself, however, an accurate endpoint forecast does not specify whether the favourable trading opportunity occurs early, late or uniformly across the execution horizon. The distinction matters whenever a prediction system is evaluated using a statistical target that contains less information than the decision requires.

Optimal-execution research has long recognised the relevance of predictive information. Trading controls can respond to expected order flow, persistent signals and price changes over the remaining horizon \citep{cartea2016,lehalle2019,forde2022,abijaber2025}. A separate literature on decision-focused prediction shows why low forecast error need not imply a low downstream decision loss \citep{elmachtoub2022,mandi2024}. These are established motivations, rather than gaps discovered by this study. The more specific problem addressed here is how to make the omitted timing information visible in a small, reproducible evaluation while keeping terminal forecasts exactly comparable.

We approach that problem through a matched-endpoint construction. Starting from a learned vector of expected execution-slot prices, we form three additional profiles with the same final coordinate. A parallel profile assigns the endpoint value to every slot; a ramp imposes a linear path to that endpoint; and a mirror reverses the intermediate contrasts around the same endpoint. Consequently, all profiles receive identical terminal-return errors and directional-hit indicators on every observation. They can nevertheless induce different feasible schedules because their relative expected prices differ. The profiles are controlled diagnostic constructions, not a set of equally plausible trained forecasting competitors.

Matching endpoints identifies what an endpoint metric cannot distinguish, but does not determine whether acting on the learned profile is worthwhile. For that second question, we decompose the realised gain from a schedule change into a timing-price payoff, a first-order baseline-cost penalty and a quadratic deviation cost. The decomposition gives an explicit economic condition: full exposure is beneficial only when the timing payoff exceeds both cost terms. Positive gross alignment with realised prices is insufficient. This framework also separates a monetary execution charge from an inventory preference, so an improvement in the combined objective is not automatically a cash saving.

The empirical illustration uses public one-minute Binance observations for Bitcoin and Ether from January through August 2026. Six rolling windows per asset allocate one month to forecast estimation, the next to exposure calibration and the third to evaluation. The final evaluation sample contains 17,664 hypothetical order episodes over 184 calendar dates. The signal is a simple imbalance measure from five completed bars, and the schedule uses fifteen subsequent minute-open proxies. The deliberately limited forecasting model makes the comparison inspectable; its simplicity is not evidence that it is the best available cryptocurrency predictor.

Under the primary assumed impact and inventory parameters, the learned schedule earns an average gross timing payoff of 0.00370 basis points but incurs a curvature cost of 0.00805 basis points. Its net objective gain is therefore negative. A half-exposure schedule reduces both benefit and cost, while a historical plug-in coefficient selects positive exposure in six of twelve asset--month windows. The more conservative lower-percentile coefficient selects zero throughout. These outcomes support a reporting lesson about the cost of acting on a forecast; they do not establish superiority of a new trading rule. In particular, a positive realised timing payoff is an ex-post sample quantity, not proof of stable exploitable predictive information.

The study contributes an operational diagnostic, a worked benefit--cost accounting and a transparent multi-window illustration. Its main value is to make the evaluation target explicit: a required-order timing decision depends on temporal price contrasts and the costs induced by the chosen response. The mathematics does not constitute a new optimal-execution principle, and the crypto application itself has direct precedents \citep{bundi2024,schnaubelt2022}. The evidence is also retrospective: the expanded design followed an initial inspected Bitcoin window. Complete results, including adverse and inconclusive comparisons, are therefore reported as exploratory rather than presented as independent confirmation.

The rest of the paper follows the decision from context to evidence. Section~\ref{sec:literature} positions the diagnostic within predictive execution, decision-focused learning and crypto microstructure. Section~\ref{sec:methods} develops the fixed-volume accounting and the chronological replay. Section~\ref{sec:results} presents the matched-profile outcomes, variation across windows, exposure policies and cost sensitivity, followed by their economic interpretation. Section~\ref{sec:conclusion} states what the evidence establishes and what remains to be tested.
